\section{Related Work}
\label{sec:related_work}

\subsection{CNC Process Monitoring and Attack Detection}
\label{sec:rw_cnc_anomaly}

CNC process monitoring has traditionally focused on tool wear, chatter, and surface quality prediction using vibration, acoustic emission, and cutting force signals~\cite{teti2010advanced,dimla2000sensor}.
Cuka and Kim~\cite{cuka2022fuzzy} demonstrated fuzzy logic-based tool condition monitoring via multi-sensor fusion, and Wang et al.~\cite{wang2023deep} surveyed the shift from hand-crafted statistical features to deep learning representations for process monitoring.

More recently, researchers have addressed \emph{adversarial} threats to CNC systems.
Energy-based methods monitor spindle power to detect unauthorized G-code modifications~\cite{kimmell2025}, achieving strong detection of gross toolpath changes but missing subtle coordinate perturbations that preserve the power envelope.
CNN-based vibration analysis~\cite{coelho2024} detects tool breakage and collisions but requires retraining on labeled anomaly examples for each new attack type.
The broader manufacturing cybersecurity literature has documented the threat landscape: Yampolskiy et al.~\cite{yampolskiy2018security} taxonomized additive manufacturing attacks by surface and impact; Belikovetsky et al.~\cite{belikovetsky2017dr0wned} demonstrated a complete attack chain inducing fatigue failure while maintaining external conformance; Graves et al.~\cite{graves2023cyberphysical} identified G-code integrity verification as an open problem.

The critical gap is the distinction between \emph{fault detection} and \emph{integrity verification}.
Process monitoring detects deviations from expected statistical behavior but does not verify that what was executed matches what was \emph{commanded}.
An adversary producing a slightly different but physically plausible toolpath may evade fault detectors entirely, because the process appears ``normal''---it is simply the wrong program.
Our reconstruction-mismatch approach addresses this gap by learning the inverse mapping from sensor signals to G-code, detecting any discrepancy between commanded and executed programs.

\subsection{Manufacturing Side-Channel Analysis}
\label{sec:rw_sidechannel}

Offensive side-channel research has shown that CNC and additive manufacturing processes leak substantial information through acoustic, electromagnetic, and power channels.
Acoustic emissions can reconstruct 3D printer toolpath geometry with sub-millimeter accuracy~\cite{al2016acoustic}; even a nearby smartphone captures sufficient data to reconstruct printed geometry~\cite{hojjati2016leave}.
Actuator power signatures can detect sabotage in additive manufacturing~\cite{gatlin2019detecting}, and spindle power profiles discriminate gross toolpath changes but not small coordinate perturbations~\cite{moore2017power}.

Our work \emph{inverts} this paradigm: rather than extracting proprietary information from sensor signals (IP theft), we use sensor signals to \emph{verify} G-code integrity (defensive monitoring).
This leverages the same physical coupling between machining parameters and sensor responses, applied for security rather than espionage.
The closest analogue is network intrusion detection via language-model perplexity on protocol sequences~\cite{brown2018recurrent}---we apply the same principle to G-code tokens reconstructed from physical sensor measurements.

\subsection{Sequence-Level Anomaly Detection}
\label{sec:rw_sequence_anomaly}

Reconstruction-error-based anomaly detection is well established~\cite{chalapathy2019deep}: autoencoders threshold reconstruction loss to identify out-of-distribution inputs~\cite{an2015variational}, while LSTMs with dynamic thresholding capture temporal dependencies in spacecraft telemetry~\cite{hundman2018detecting}, and Transformer architectures such as TranAD use self-conditioning and adversarial training for multivariate time-series anomaly detection~\cite{tuli2022tranad}.
In the discrete domain, language models score anomalies via perplexity or NLL, flagging ``surprising'' sequences~\cite{brown2018recurrent}, an approach successfully applied to system log anomaly detection.

Our reconstruction target differs from both continuous time series and free-form text: G-code is a \emph{structured token sequence} with strict grammar.
Each line begins with a command token (G0, G1, etc.) followed by parameter--value pairs, where parameter letters and numeric values occupy distinct fields.
This structure provides an additional grammar-violation detection signal unavailable in other domains, and motivates our multi-head decoder with separate heads for token type, command, parameter, and numeric value.
No prior work has applied autoregressive token-level anomaly scoring to G-code sequences reconstructed from multi-modal sensor inputs; the closest analog uses language models for IT log-file anomaly detection~\cite{brown2018recurrent}, which lacks G-code's mixed vocabulary and grammatical constraints.

\subsection{Calibration and Confidence Estimation}
\label{sec:rw_calibration}

Modern neural networks are often poorly calibrated: softmax probabilities do not reflect true predictive accuracy~\cite{guo2017calibration}.
Temperature scaling---a single learned scalar on a validation set---is widely recommended as a default calibration strategy~\cite{guo2017calibration,platt1999probabilistic,naeini2015obtaining}.
For anomaly detection, calibration matters because thresholds are set from score distributions; miscalibrated scores make operating-point selection unreliable.
OOD detection methods based on softmax thresholding~\cite{hendrycks2017baseline}, Mahalanobis distance~\cite{lee2018simple}, ODIN perturbation~\cite{liang2018enhancing}, and energy scoring~\cite{liu2020energy} all depend on well-calibrated confidence estimates.

As we show in Section~\ref{sec:results_calibration}, temperature scaling produces mixed results on our decoder, cautioning against uncritical application of post-hoc calibration methods.
